% ==========================================================================
% ReasonForge: Verifiable AI Reasoning with Zero-Knowledge Attestation
% Workshop paper (4 pages) -- target: AISec @ CCS / NeurIPS workshops
% ==========================================================================
\documentclass[10pt,a4paper,twocolumn]{article}

% ── Packages ──────────────────────────────────────────────────────────────
\usepackage[utf8]{inputenc}
\usepackage[T1]{fontenc}
\usepackage{lmodern}
\usepackage{amsmath,amssymb}
\usepackage{booktabs}
\usepackage{hyperref}
\usepackage{tikz}
\usetikzlibrary{arrows.meta,positioning}
\usepackage[margin=0.75in]{geometry}
\usepackage{xcolor}
\usepackage{enumitem}
\usepackage{microtype}
\usepackage{natbib}
\bibliographystyle{plainnat}

\definecolor{rfblue}{HTML}{2563EB}
\definecolor{rfdark}{HTML}{1E293B}
\definecolor{rfgreen}{HTML}{059669}
\definecolor{rfred}{HTML}{DC2626}

% ── Title ─────────────────────────────────────────────────────────────────
\title{\textbf{ReasonForge: Verifiable AI Reasoning\\
with Zero-Knowledge Attestation}}

\author{
  Ravi Shankar Bejini\\
  \texttt{ravishankarbejini@gmail.com}\\
  \small\url{https://github.com/UNSPECIFIEDCODER/ReasonForge}
}

\date{March 2026}

% ==========================================================================
\begin{document}
\maketitle

% ── Abstract ──────────────────────────────────────────────────────────────
\begin{abstract}
Large language models provide no formal guarantees about the validity of
their intermediate reasoning steps.  We present ReasonForge, a system for
mechanically verifying multi-step AI reasoning chains.  ReasonForge
translates reasoning outputs into domain-appropriate formal
representations---Lean~4 proofs, executable tests, and SMT
formulas---and produces verification certificates encapsulated in
Groth16 zero-knowledge proofs.  We describe a production code
verification pipeline with AST-based security scanning, automated test
generation, and sandboxed execution, and present an implementation
comprising 7{,}418 lines of Python, a 153-line circom ZK circuit, and
316 automated tests.
\end{abstract}

% ======================================================================
\section{Introduction}
% ======================================================================

Current evaluation paradigms assess LLM outputs at the level of final
answers, not intermediate reasoning steps.  A 20-step reasoning chain
with 95\% per-step accuracy yields only 35.8\% end-to-end reliability.
Process reward models~\citep{lightman2023verify} improve step-level
evaluation but remain statistical estimators that cannot produce
externally auditable attestations.

ReasonForge treats AI reasoning verification as an infrastructure problem
with a cryptographic solution.  The system (i)~translates structured
reasoning chains into mechanically verifiable representations across
three backends (Lean~4, executable tests, SMT solvers), (ii)~implements a
production code verification pipeline, and (iii)~produces Groth16 ZK
proofs over verification certificates for privacy-preserving attestation.

% ======================================================================
\section{System Architecture}
% ======================================================================

\subsection{Verification Pipeline}

ReasonForge accepts structured reasoning traces with per-step domain
classification (mathematical, code, or logical).  Each step is routed to
the appropriate verification backend:

\begin{itemize}[nosep,leftmargin=*]
  \item \textbf{Mathematics $\rightarrow$ Lean~4:} Proof terms are
    type-checked; successful checking constitutes formal proof.
  \item \textbf{Code $\rightarrow$ Sandbox + Tests:} Source code is
    scanned for security vulnerabilities, tested via auto-generated
    pytest suites, and executed in an isolated subprocess.
  \item \textbf{Logic $\rightarrow$ SMT:} First-order formulas are
    submitted to Z3/CVC5; validity is confirmed via UNSAT on the
    negation.
\end{itemize}

\subsection{Code Verification Pipeline}

The production instantiation targets AI-generated Python code with four
stages:

\paragraph{Security Scanner.} An AST walker applies 10 detection rules
(RF-SEC-001 through RF-SEC-010) covering \texttt{eval}/\texttt{exec}
calls, dangerous imports, \texttt{pickle.loads}, hardcoded credentials,
and bare \texttt{except} clauses.  Each vulnerability is tagged with
severity (critical/high/medium/low) and line number.

\paragraph{Test Generator.} Static analysis extracts function signatures
and produces five categories of pytest tests: callable, type-check,
None-input, docstring-example, and edge-case.  No LLM is required.

\paragraph{Sandbox Executor.} A subprocess isolates execution with
meta-path import hooks blocking 10 dangerous modules and wall-clock
timeout enforcement.

\paragraph{Confidence Scoring.} The composite score
$C = 0.2 S + 0.3\,\text{Sec} + 0.5 T$ weights syntax validity~($S$),
security score~(Sec), and test pass rate~($T$).  Verdict is
\textsc{passed} when $S=1$, risk $\notin \{\text{critical, high}\}$, and
$T \geq 0.8$.

\subsection{ZK Attestation}

A 153-line circom circuit generates Groth16 proofs over certificate
fields.  Public inputs include the task hash (8$\times$32-bit limbs),
overall verdict, total/verified step counts, and timestamp.  Private
inputs are the 20-element binary step verdict array.  The circuit
enforces: (i)~binary constraints on step verdicts, (ii)~sum equals
verified count, (iii)~verdict consistency, (iv)~range checks on hash
limbs, and (v)~indicator exclusivity.  The proof attests ``\emph{trace
with hash $H$ received verdict $V$ with $k$/$n$ steps verified}''
without revealing step details.

% ======================================================================
\section{Incentive Mechanism}
% ======================================================================

The decentralised network scores miners via the Composite Miner Score:

\vspace{-4pt}
\begin{equation}
  \text{CMS}(m,t) = 0.45 Q + 0.30 A + 0.15 N + 0.10\,\text{Eff}
\end{equation}
\vspace{-4pt}

\noindent and the Composite Translation Score for formalisation quality:

\vspace{-4pt}
\begin{equation}
  \text{CTS}(m,t) = 0.40 C_{\text{comp}} + 0.25 C_{\text{corr}}
    + 0.20 C_{\text{cmpl}} + 0.15 C_{\text{eff}}
\end{equation}
\vspace{-4pt}

Adversarial resistance is provided by 15\% trap problem injection,
blind evaluation via cryptographic identity hashing, quadratic validator
slashing ($\text{slash} \propto (\theta - \text{VAS})^2$), and a
similarity penalty for plagiarised submissions.

% ======================================================================
\section{Evaluation}
% ======================================================================

The implementation comprises 7{,}418 lines of Python (37 modules), a
153-line circom circuit, and 316 tests (310 passing, 6 requiring
circom/snarkjs).  The full test suite executes in $\sim$12 seconds on
commodity hardware.

\begin{table}[h]
\centering
\small
\begin{tabular}{@{}lr@{}}
\toprule
\textbf{Component} & \textbf{Tests} \\
\midrule
Certificates + ZK   & 77 \\
Code verify (unit)   & 38 \\
Code verify (integ.) & 7 \\
Engine + Simulator   & 50 \\
Protocol + Types     & 45 \\
Translation          & 32 \\
Verification         & 43 \\
Proof layer sim.     & 24 \\
\midrule
\textbf{Total}       & 316 \\
\bottomrule
\end{tabular}
\caption{Test suite breakdown by component.}
\end{table}

We evaluated the code verification pipeline on the full OpenAI
HumanEval benchmark~\citep{humaneval2021} (164 canonical Python
solutions).  All 164 functions pass syntax validation (100\%) and 161
(98.2\%) are classified as security-safe.  55 functions (33.5\%)
receive a \textsc{passed} verdict; the remaining 109 (66.5\%) are
classified as \textsc{failed} due to auto-generated robustness tests
(None-input, edge-case, type-check) that the canonical solutions do
not handle.  This demonstrates that the pipeline measures
\emph{robustness} beyond correctness: functions that pass verification
are defensively programmed.  Mean confidence is 0.817; median latency
is 2.2\,s per function.

The certificate round-trip (pipeline $\rightarrow$ report $\rightarrow$
ZK certificate $\rightarrow$ registry $\rightarrow$ verification)
completes end-to-end for all test scenarios.

% ======================================================================
\section{Conclusion}
% ======================================================================

ReasonForge demonstrates that mechanical verification of AI reasoning
chains is practical and deployable.  The code verification pipeline
provides immediate value for AI code generation workflows, while the ZK
attestation layer enables privacy-preserving compliance for enterprise
deployments.  Future work includes improving translation fidelity for
advanced mathematics, extending code verification to additional
languages, and migrating to transparent proof systems to eliminate the
trusted setup.

\smallskip
\noindent\textbf{Code:}
\url{https://github.com/UNSPECIFIEDCODER/ReasonForge}

% ── References ────────────────────────────────────────────────────────────
\bibliography{refs}

\end{document}
